\chapter{Introducción}
\label{cap:introduccion}

\chapterquote{Frase célebre dicha por alguien inteligente}{Autor}

Según la normativa  para Trabajos de Fin de Grado\footnote{\url{https://informatica.ucm.es/file/normativatfg-2021-2022?ver} (ver versión actualizada para cada curso académico)}, la memoria incluirá una portada normalizada con la siguiente información: título en castellano, título en inglés, autores, profesor director, codirector si es el caso, curso
académico e identificación de la asignatura (Trabajo de fin de grado del Grado en -
nombre del grado correspondiente-, Facultad de Informática, Universidad Complutense de Madrid). Los datos referentes al título y director (y codirector en su caso) deben corresponder a los publicados en la página web de TFG.

La memoria debe incluir la descripción detallada de la propuesta hardware/software realizada y ha de contener:
\begin{enumerate}
\renewcommand{\theenumi}{\alph{enumi}}
	\item un índice,
	\item un resumen y una lista de no más de 10 palabras clave para su búsqueda bibliográfica, ambos en castellano e inglés,
	\item una introducción con los antecedentes, objetivos y plan de trabajo,
	\item resultados y discusión crítica y razonada de los mismos, con sus conclusiones,
	\item bibliografía.
\end{enumerate}
Para facilitar la escritura de la memoria siguiendo esta estructura, el estudiante podrá usar las plantillas en LaTeX o Word preparadas al efecto y publicadas en la página web de TFG.

La memoria constará de un mínimo de 25 páginas para los proyectos realizados por un único estudiante, y de al menos 5 páginas más por cada integrante adicional del grupo. En este número de páginas solo se tiene en cuenta el contenido correspondiente a los apartados c y d del punto anterior.

La memoria puede estar escrita en castellano o inglés, pero en el primer caso la introducción y las conclusiones deben aparecer también en inglés. Las memorias de los TFG matriculados en el grupo I deberán estar escritas íntegramente en inglés, excepto por lo especificado en los puntos 1 y 2 anteriores (título, resumen y lista de palabras clave).

En caso de trabajos no unipersonales, cada participante indicará en la memoria su contribución al proyecto con una extensión de al menos dos páginas por cada uno de los participantes.

Todo el material no original, ya sea texto o figuras, deberá ser convenientemente citado y referenciado. En el caso de material complementario se deben respetar las licencias y \emph{copyrights} asociados al software y hardware que se emplee. En caso contrario no se autorizará la defensa, sin menoscabo de otras acciones que correspondan.


\section{Motivación}
Al plantear este Trabajo de Fin de Grado, existía una premisa clara: que el resultado tuviera una aplicación real. No se trataba de resolver un problema inventado para cumplir un requisito académico, sino de utilizar los conocimientos adquiridos durante la carrera para abordar una necesidad concreta de alguien cercano. La idea era que el trabajo sirviera como ejercicio técnico y, al mismo tiempo, como solución a un problema que alguien experimentaba en su día a día profesional.

Ese alguien resultó ser un amigo que trabaja como nutricionista. En una conversación sobre su rutina diaria, quedó en evidencia un problema que él mismo no percibía como tal, porque lo había normalizado: su flujo de trabajo estaba fragmentado entre cinco o más herramientas que no se comunicaban entre sí. Utilizaba formularios de Google para recoger información de sus clientes, Google Calendar para gestionar las citas, WhatsApp como canal de comunicación directa, un cuaderno físico para sus notas privadas sobre cada caso, y la plataforma de gestión de su clínica para el resto. Para hacerse una idea completa de la situación de un solo cliente, tenía que consultar todas estas fuentes manualmente y reconstruir la información por su cuenta. Además, una parte significativa de su tiempo se iba en responder consultas repetitivas ---qué comprar en el supermercado, si se podía sustituir un ingrediente, cuándo era la próxima cita--- que podrían resolverse sin su intervención directa.

A partir de esa conversación surgió la idea central de este proyecto: diseñar una plataforma que centralizara todas esas herramientas dispersas en un único punto de acceso. Si toda la información de los clientes ---formularios, citas, mensajes, notas, planes nutricionales, registros diarios--- viviera en una misma base de datos, no solo se eliminaría la fragmentación, sino que se abrirían posibilidades imposibles con herramientas desconectadas: detectar patrones de forma automática, generar alertas cuando un cliente muestra señales de abandono, o construir un perfil completo sin tener que cruzar datos a mano. En definitiva, que el nutricionista pudiera dedicar su tiempo a lo que realmente importa ---diseñar dietas, acompañar a sus clientes, lograr que los tratamientos funcionen--- en lugar de invertirlo en tareas administrativas que una aplicación bien diseñada puede resolver.


\section{Objetivos}
Este Trabajo de Fin de Grado se desarrolla en el marco del Doble Grado en Ingeniería Informática y Administración y Dirección de Empresas, lo cual permite abordarlo desde dos perspectivas complementarias: la técnica, apoyada en los conocimientos adquiridos en la asignatura de Ampliación de Bases de Datos, y la de negocio, aplicando herramientas de análisis de mercado y viabilidad de producto propias de la formación en ADE. El objetivo no es simplemente cumplir un requisito académico, sino aprovechar esta doble formación para construir un proyecto con aplicación real.

El objetivo general de este trabajo es diseñar e implementar la base de datos de una aplicación web orientada a la gestión integral de consultas de nutrición, que centralice el flujo de trabajo del nutricionista profesional y mejore la experiencia del cliente. Se busca llevar a la práctica, en un proyecto desarrollado desde cero, los conocimientos técnicos de diseño y gestión de bases de datos junto con la capacidad de análisis de mercado y enfoque de producto, de modo que el resultado no solo sea técnicamente sólido sino también viable como herramienta real.

Para alcanzar este objetivo general, se plantean los siguientes objetivos específicos:

\begin{enumerate}
    \item Analizar el flujo de trabajo real de un nutricionista profesional para extraer los requisitos funcionales del sistema. Esto implica comprender cómo trabaja en su día a día, qué herramientas utiliza actualmente, qué problemas enfrenta y qué necesita realmente. Los requisitos no se inventan en abstracto, sino que se derivan de la observación directa de un caso real, lo que garantiza que el diseño responda a necesidades concretas y no a suposiciones.

    \item Realizar un estudio de mercado y análisis de la competencia en el ámbito de las plataformas de nutrición, aplicando conocimientos propios de la formación en ADE. Esto incluye identificar los principales competidores, analizar sus fortalezas y debilidades, detectar carencias y oportunidades en el mercado, y posicionar NutriApp con funcionalidades diferenciadoras que aporten valor real. El objetivo es que la propuesta no solo sea técnicamente correcta, sino que tenga sentido como producto viable en un mercado existente.

    \item Diseñar un esquema de base de datos relacional completo que refleje toda la lógica de negocio de la aplicación. Esto abarca el modelado de las entidades y sus relaciones, el uso de triggers para automatizar comportamientos ---como la generación de alertas a partir de respuestas de formularios---, consultas complejas con agregaciones temporales ---como el cálculo de rachas de check-ins consecutivos---, y relaciones avanzadas como las muchos-a-muchos auto-referenciadas del sistema de sustituciones de alimentos. El diseño de la base de datos constituye el eje central del proyecto y debe demostrar un dominio profundo de los conceptos trabajados en Ampliación de Bases de Datos.
\end{enumerate}


\section{Plan de trabajo}
El desarrollo de este proyecto ha seguido un proceso iterativo en el que cada fase se apoya en los resultados de la anterior. A continuación se describe el trabajo realizado hasta el momento, organizado en las fases que han guiado el avance del proyecto.

\subsection*{Fase 1 --- Investigación con nutricionistas reales}

El punto de partida fue entender en profundidad el problema que se quería resolver. Para ello, se mantuvo una serie de conversaciones detalladas con el nutricionista cuyo caso motivó el proyecto, con el objetivo de conocer a fondo su método de trabajo: qué herramientas utilizaba, en qué puntos perdía más tiempo, qué tareas le resultaban más ineficientes y qué información se le escapaba por tenerla repartida entre plataformas desconectadas. Cuanto mejor se comprendiera su flujo de trabajo y sus puntos débiles, más preciso sería el diseño de la solución.

Sin embargo, basar el producto únicamente en la experiencia de una persona suponía un riesgo: las necesidades de un nutricionista concreto no tienen por qué ser representativas del colectivo. Para ampliar la perspectiva, se contactó a través de LinkedIn con otros profesionales del sector, a quienes se les preguntó qué aspectos de su día a día podrían mejorar y qué funcionalidades valorarían en una herramienta de gestión integral. Esta ronda de consultas permitió contrastar las observaciones iniciales, incorporar necesidades que no habían surgido en el primer caso y construir una visión del producto mucho más amplia y fundamentada.

\subsection*{Fase 2 --- Análisis de competencia y mercado}

Una vez definidas las necesidades desde el lado del usuario, el siguiente paso fue investigar el estado actual del mercado. Era razonable asumir que ya existían soluciones orientadas a nutricionistas, y antes de diseñar nada convenía saber quiénes eran los competidores, qué ofrecían y dónde dejaban huecos sin cubrir.

Se analizaron las principales plataformas del sector ---Nutrium, Indya, Dietopro, Nutrimeal y Practice Better, entre otras--- estudiando sus funcionalidades, sus modelos de negocio, sus puntos fuertes y sus carencias. El objetivo de este análisis no era copiar lo que otros hacían bien ni competir directamente con ellos, sino detectar qué aspectos se habían quedado fuera del radar durante la fase de investigación con nutricionistas e incorporarlos a la propuesta. Los competidores sirvieron, sobre todo, como contraste: si un problema identificado en las entrevistas no estaba resuelto por ninguna plataforma existente, eso confirmaba que había una oportunidad real. Al final de esta fase, la idea del producto estaba mucho más definida y fundamentada, con una propuesta de valor clara y diferenciada.

\subsection*{Fase 3 --- Diseño visual de la aplicación}

Con los requisitos funcionales extraídos y el posicionamiento del producto definido, se pasó a traducir todo ello en una interfaz concreta. Se diseñaron los mockups de la aplicación: las pantallas que reflejan todas las herramientas e interfaces que tanto el nutricionista como el cliente tendrían a su disposición. El objetivo era doble. Por un lado, que el nutricionista pudiera gestionar todo su trabajo desde un único punto de acceso ---clientes, citas, formularios, planes, mensajes y notas--- sin necesidad de recurrir a herramientas externas. Por otro, que el cliente contara con herramientas propias que le permitieran consultar su plan, registrar sus hábitos, hacer seguimiento de su progreso y resolver dudas frecuentes sin tener que contactar al nutricionista cada vez.

Un aspecto clave en esta fase fue garantizar que el flujo de información entre ambos paneles fuera lo más natural posible: que los datos introducidos por el cliente llegaran al nutricionista de forma organizada, y que las indicaciones del nutricionista se reflejaran de forma clara en la interfaz del cliente. Cada pantalla se diseñó para cubrir los casos de uso identificados en las fases anteriores, de modo que el diseño visual no fuera un ejercicio estético sino una traducción directa de los requisitos funcionales.

\subsection*{Fase 4 --- Ampliación y refinamiento del diseño}

Una vez completado el primer conjunto de pantallas, se identificó que varias secciones de la aplicación ---planes nutricionales, informes, ajustes del nutricionista, perfil del cliente, citas del cliente y formularios del cliente--- no estaban representadas en el prototipo inicial. Se diseñaron las pantallas correspondientes siguiendo el mismo sistema de diseño (tipografía, colores, componentes y disposición) que las pantallas originales, garantizando la coherencia visual del conjunto.

Además, se realizó una revisión iterativa del diseño existente para incorporar mejoras detectadas durante el análisis de requisitos. Entre los cambios más relevantes se encuentran: la incorporación de vistas diaria, semanal y mensual en el calendario del nutricionista junto con la posibilidad de añadir citas directamente; la adición de un botón de envío de formularios tanto en la biblioteca de formularios como en la lista de clientes; la reestructuración de la pantalla de planes nutricionales en dos pestañas (planes de dieta completos y plantillas de comidas individuales); y la transformación de la vista del plan del cliente de una vista exclusivamente semanal a una vista diaria por defecto con opción de cambiar a semanal o mensual. También se sustituyó la opción de cierre de sesión del panel del cliente por una pantalla de ajustes donde el cliente puede gestionar su cuenta, cambiar su contraseña y configurar sus preferencias.

El resultado de esta fase fue un primer prototipo navegable con el conjunto básico de pantallas interconectadas ---panel del nutricionista y panel del cliente--- en el que todos los requisitos funcionales definidos inicialmente en el capítulo~\ref{cap:analisisRequisitos} estaban representados en al menos una pantalla del diseño.

\subsection*{Fase 5 --- Segunda iteración: análisis de competencia en profundidad y nuevas funcionalidades}

Tras completar el primer prototipo, se llevó a cabo un segundo análisis de competencia, esta vez mucho más profundo y contrastando NutriApp con siete plataformas de referencia (Nutrium, Practice Better, Healthie, INDYA, Dietopro, Cronometer Pro y Foodzilla). El objetivo no era volver a situar el producto en el mercado, sino revisarlo a la luz de lo que ya existe y detectar funcionalidades que merecía la pena incorporar porque aportaban valor real al nutricionista o al cliente, incluso aunque no estuvieran en la propuesta inicial.

De este análisis surgieron cuatro incorporaciones importantes al diseño, que se añadieron siguiendo estrictamente el sistema de diseño existente para mantener la coherencia visual. La primera fue un \textbf{motor de automatizaciones} configurable desde los ajustes del nutricionista, inspirado en Practice Better: permite definir reglas del tipo \textit{cuando ocurre X, hacer Y} para modelar procedimientos recurrentes de la práctica (bienvenida a nuevos clientes, seguimiento tras sesión, detección de baja adherencia sostenida). La segunda fue un \textbf{marketplace de recursos educativos}, inspirado en Healthie, donde el nutricionista puede publicar sus propios materiales formativos y gestionar sus ventas, y el cliente dispone de una sección destacada con los recursos publicados específicamente por su profesional. La tercera fue una \textbf{base de datos de alimentos con reconocimiento inteligente}, integrada con Open Food Facts y complementada con dos vías de captura rápida ---escaneo de códigos de barras y reconocimiento de alimentos por fotografía mediante un modelo de inteligencia artificial---, accesible tanto al nutricionista (para añadir alimentos personalizados a sus planes) como al cliente (para registrar y comparar alimentos con sus objetivos nutricionales del día). La cuarta fue la incorporación de \textbf{videollamadas y llamadas de voz integradas en el chat}, que evita la dispersión en herramientas externas como Zoom o Google Meet.

Además de estas cuatro innovaciones, la iteración incluyó otros refinamientos relevantes: una nueva pantalla de \textbf{ficha detallada del cliente} en el panel del nutricionista, que consolida en una sola vista las notas privadas, el plan actual, la próxima sesión con su foco, las alertas recientes y el histórico de check-ins; un canal de \textbf{comunidad entre clientes del mismo nutricionista} en la mensajería, pensado como apoyo entre pares y con aislamiento garantizado entre consultas de profesionales distintos; un botón de \textit{swap} en cada comida del plan diario del cliente, que le permite cambiarla por una opción equivalente preconfigurada por el nutricionista; y una reestructuración completa del dashboard del cliente para dar una jerarquía visual más clara a la información relevante del día (check-in, próxima cita, comidas planificadas y racha de adherencia).

El resultado de esta fase es un prototipo navegable completo con veinticuatro pantallas interconectadas ---once para el nutricionista, doce para el cliente y una pantalla de entrada--- en el que los sesenta y un requisitos funcionales definidos en el capítulo~\ref{cap:analisisRequisitos}, organizados en catorce bloques temáticos, están representados al menos en una pantalla del diseño. El prototipo incluye también toda la navegación entre pantallas cableada como reacciones de Figma, de forma que un evaluador externo puede recorrerlo como si fuera una aplicación real.

\subsection*{Fase 6 --- Validación con usuarios reales}

Con el prototipo completo, se preparó un proceso de validación con usuarios reales para recoger feedback sobre la intuición de los flujos y la calidad del diseño. Para ello se redactaron dos guías de prueba de usabilidad ---una para el panel del nutricionista y otra para el panel del cliente--- en formato editable (\texttt{.docx}), cada una con una docena de flujos de navegación planteados en orden aleatorio. Los flujos se formulan como pequeños retos (``quieres consultar la adherencia media de tus clientes, ¿a qué sección del menú lateral irías?''), de modo que el evaluador tenga que tomar una decisión de navegación en cada caso y marcar si le ha resultado obvia, si ha tenido que pasar por una pantalla equivocada antes, o si directamente no sabría cómo llegar. Ambos documentos se preparan en castellano e inglés para poder enviarse a profesionales de distintos contextos.

%% PENDIENTE: añadir fases restantes (diseño de base de datos, implementación, pruebas, etc.) %%